% Complete local proof dependencies for the selected IEL continuation.
\section{Question-to-question implication: proofs}
\label{sec:p18-erotetic-implication-annex}

Use the fixed language, direct-answer assignment and admissible partition class
of Section~\ref{sec:p18-question-to-question-implication}. Entailment is the
partition-based relation of Definition~\ref{def:p18-ei-setting}. The first
source supplies the native implication/regularity concepts and cases; the
second supplies the declarative-basis and conditional/conjunctive constructions
\cite[secs.~4,7--8]{Wisniewski2001QuestionsInferences}
\cite[secs.~2--5]{Wisniewski2018DeductionReduction}.

\begin{annexitem}{annex:p18.ei.basis}{thm:p18-ei-basis}{The source's semantic representation of consequence}
\begin{proof}
Assume \(X\vdash_\ell A\), and let an admissible partition have true
component \(T=\operatorname{Cn}_\ell(Y)\) with \(X\subseteq T\).
Premise monotonicity gives \(T\vdash_\ell A\), hence
\(Y\cup T\vdash_\ell A\). By definition \(Y\vdash_\ell B\) for every
\(B\in T\). Cut for sets gives \(Y\vdash_\ell A\), so \(A\in T\).
This proves semantic entailment for every admissible partition.

Conversely suppose \(X\not\vdash_\ell A\). Lb provides a maximally
consistent \(Y\supseteq X\) with \(Y\not\vdash_\ell A\). Overlap gives
\(Y\subseteq\operatorname{Cn}_\ell(Y)\), so the partition generated by
\(Y\) satisfies \(X\) but not \(A\). It refutes semantic entailment.
The proof uses the stipulated Lb extension, not a purported existence proof
for maximal consistent sets in every logic.
\end{proof}
\end{annexitem}

\begin{annexitem}{annex:p18.ei.profile}{thm:p18-ei-profile}{Soundness, witness force and edge cases}
\begin{proof}
Fix a partition with \(X\subseteq T\). If \(Q\) is sound, select a true
\(A\in dQ\). The first implication clause supplies a true member of
\(dR\). Conversely, if \(R\) is sound, select a true \(B\in dR\).
Its nonempty proper witness contains a true original answer, so \(Q\) is
sound. These arguments also prove that the two questions are unsound together
when \(X\) is true.

For regular implication, choose the singleton containing the entailed
original answer. It is nonempty and proper since \(|dQ|\geq2\). Adding
premises only restricts the partitions to be checked, so every entailment and
every existing witness is preserved. If no partition satisfies \(X\), all
these entailments hold vacuously; fix any original answer to obtain the required
singleton witnesses. Evocation fails there because every direct answer is
entailed, contradicting its non-entailment clause.

Assume evocation of \(Q\) and regular implication to \(R\). Coverage of
\(Q\) and the first implication clause give coverage of \(R\). If
\(X\models_L B\) for a direct answer of \(R\), its regular witness gives
\(X,B\models_L A\) for some \(A\in dQ\). At every partition satisfying
\(X\), both \(B\) and then \(A\) are true. Thus \(X\models_L A\),
contrary to evocation. This proves evocation of \(R\).

For failure with general implication, use the displayed
\(Q=?\{p,q,r\}\), \(R=?\{p\lor q,r\}\), \(X=\{p\lor q\}\).
Each original answer implies a target answer; the target witnesses are
\(\{p,q\}\) and \(\{r\}\). Coverage of \(Q\) holds. Assignments
\((p,q,r)=(0,1,0)\) and \((1,0,0)\) falsify respectively \(p\) and
\(q\), and either falsifies \(r\), while satisfying \(X\). Hence \(Q\)
is evoked. But \(p\lor q\) is entailed by \(X\), so \(R\) is not.
At \((1,1,1)\), the first target answer and the original answer outside its
witness are simultaneously true. This proves that witness reduction is not
general negative elimination.

Finally, the empty conclusion-set intersects no true component. Its
multiple-conclusion entailment therefore holds precisely when there is no
partition satisfying its premises. Such a set cannot be an allowed witness
or a question in this language. A well-formed question with all answers
false simply lacks soundness at that partition; its syntax is unaffected.
\end{proof}
\end{annexitem}

\begin{annexitem}{annex:p18.ei.presuppositions}{thm:p18-ei-presuppositions}{The exact presupposition characterizations}
\begin{proof}
If \(X\) evokes \(Q\), every \(X\)-partition contains an original
answer, and that answer entails each presupposition. Thus \(X\) entails
each presupposition. If \(Q\) is normal and \(X\) entails all its
presuppositions, every \(X\)-partition satisfies \(\operatorname{Pres}(Q)\)
and hence some direct answer. With the stipulated individual non-entailments,
this gives evocation. The reverse was just proved. If \(Q\) is regular,
evocation entails its prospective presupposition; conversely entailing one
such presupposition gives coverage, and the same non-entailments give evocation.

For normal \(Q,R\), assume the first implication clause. At any partition
satisfying \(X\cup\operatorname{Pres}(Q)\), normality gives an answer
\(A\) to \(Q\). The first clause gives a true answer to \(R\), which
entails every \(C\in\operatorname{Pres}(R)\). Thus all the displayed
presupposition entailments hold. Conversely, at any partition satisfying
\(X,A\) for \(A\in dQ\), every presupposition of \(Q\) is true. The
displayed entailments make every presupposition of \(R\) true. Normality
of \(R\) gives a direct answer, proving the first implication clause.

For regular questions, choose prospective presuppositions \(D\) and \(C\).
Under the first clause, an \(X,D\)-partition has an answer to \(Q\), then
an answer to \(R\), then \(C\). Hence \(X,D\models_L C\). Conversely,
if prospective \(D,C\) satisfy that entailment, any \(X,A\)-partition
with \(A\in dQ\) has \(D\), then \(C\), then an answer to \(R\).
Again this proves only the first clause.

If \(dQ\) has exactly two elements, every nonempty proper witness set
for a purported partial answer is a singleton \(\{A\}\). Its two
defining entailments then say both \(C\models_L A\) and
\(A\models_L C\). That makes \(C\) just-complete, excluded by the
partial-answer definition. Thus there is no partial answer of this type.

For a concrete normal-but-not-regular case in general semantics, use Boolean
formulas over variables \(a_i,c_i\) indexed by the natural numbers, but
restrict admissible valuations to the following class. At \(W_j\), only
\(a_j\) is true among the a-variables and every c-variable is true. At
\(N_k\), all a-variables are false and only \(c_k\) is false among the
c-variables. Give a question \(Q\) the infinite direct-answer set
\(\{a_i:i\in\mathbb N\}\). Each \(c_i\) is a presupposition. Every
null valuation \(N_k\) falsifies the presupposition \(c_k\), while each
\(W_j\) satisfies all presuppositions and has an answer. Hence the full
presupposition set entails the answer set: \(Q\) is normal.

Any one presupposition \(C\) is a finite formula. Choose \(j\) outside
its finitely many a-indices and \(k\) outside its finitely many c-indices.
The valuations \(W_j,N_k\) agree on every variable in \(C\), so give it
the same truth value. It is true at \(W_j\), since \(a_j\) entails every
presupposition, and therefore true at the unsound valuation \(N_k\).
Thus no presupposition is prospective and \(Q\) is not regular. This
example uses a restricted admissible class, not all Boolean valuations,
and an infinite answer set. It is not a countercase to finite Boolean PMCE.

The independent polar countercase in the reader verifies the missing clause
explicitly. Its answer list covers every Boolean valuation. For either chosen
target answer \(s\) or \(\neg s\) and any proper subset of \(\{p,q,r\}\),
choose an excluded variable to be the only true original answer. Fix \(s\)
to make the chosen target answer true. No answer in the proposed subset is
true. Thus no reduction witness exists. Neither normality nor prospective
presupposition replaces the second clause.
\end{proof}
\end{annexitem}

\begin{annexitem}{annex:p18.ei.conditional-lemmas}{thm:p18-ei-conditional-lemmas}{All three source conditional transformations}
\begin{proof}
For nonempty \(Y\), a partition makes every member of \(F_C(Y)\)
untrue exactly when it makes \(C\) true and every member of \(Y\) untrue.
In one direction apply the conditional truth clause to any member to obtain
\(C\), and to all members to obtain their untruth; in the other apply the
clause to each member. Therefore the partitions refuting
\(X,C\mathrel{\|=}_L Y\) are exactly those refuting
\(X\mathrel{\|=}_L F_C(Y)\). This proves the first equivalence.

Any partition refuting an entailment with conclusion set \(F_C(Y)\)
already makes \(C\) true. At such a partition \(D\) is true exactly when
\(C\to D\) is true. Consequently replacing premise \(D\) by
\(C\to D\) preserves exactly the counterpartitions, proving the second
equivalence. Apply the first with premise set \(X\cup\{D\}\), then
the second, to prove the third.

Nonemptiness matters. In the Boolean case let \(X=\varnothing\),
\(C=p\land\neg p\) and \(Y=\varnothing\). Then \(\{C\}\) entails
the empty set vacuously, whereas the empty premise set does not entail it.
Since \(F_C(\varnothing)=\varnothing\), the first equivalence fails
without its stated condition.
\end{proof}
\end{annexitem}

\begin{annexitem}{annex:p18.ei.evocation-deduction}{thm:p18-ei-evocation-deduction}{General semantic evocation deduction and finite reduction}
\begin{proof}
Apply the first conditional lemma to \(Y=dQ\) to transfer coverage.
For each \(A\in dQ\), apply it to the singleton \(\{A\}\):
\(X,C\models_L A\) iff \(X\models_L C\to A\). Negating these
equivalences transfers every required non-entailment. Injectivity makes the
answers correspond one-to-one. The assumed closure supplies the transformed
question, and the inverse condition supplies an unprefixed question when
starting from a prefixed one. Hence evocation is equivalent in both directions.

For finite \(X\), remove \(C_n\) first, then \(C_{n-1}\), and so on,
each time adding its prefix to every current answer. Induction on \(n\)
gives exactly the displayed nested prefix and empty remaining basis. Each
intermediate answer set is a question by closure. Reversing the finitely many
biconditionals proves the reverse reduction. For \(n=0\), nothing changes.
There is no limiting question introduced for infinitely many premises.
\end{proof}
\end{annexitem}

\begin{annexitem}{annex:p18.ei.deduction}{thm:p18-ei-deduction}{Both directions of implication deduction and reduction}
\begin{proof}
The third conditional lemma, with \(D=A\) and \(Y=dR\), gives
\[
 X,C,A\mathrel{\|=}_L dR
 \quad\Longleftrightarrow\quad
 X,C\to A\mathrel{\|=}_L F_C(dR).
\]
As \(A\) ranges over \(dQ\), this is exactly equivalence of the two
first implication clauses.

For a target answer \(B\), suppose \(Y_B\) is an original nonempty proper
reduction witness. The third lemma transfers its entailment to
\(X,C\to B\mathrel{\|=}_L F_C(Y_B)\). The prefix map is injective, so
\(F_C(Y_B)\) is nonempty and proper in \(F_C(dQ)\). This gives the
second clause of the transformed relation.

Conversely, a witness \(W\subsetneq F_C(dQ)\) for target answer
\(C\to B\) has a unique preimage
\(Y=\{A\in dQ:C\to A\in W\}\). Its nonemptiness and properness
follow from injectivity and the equality of the entire image with
\(F_C(dQ)\). The third lemma gives
\(X,C,B\mathrel{\|=}_L Y\). Doing this for every target answer supplies
all original witnesses. Closure and inverse closure ensure the question
expressions themselves exist. Thus the two directions are proved.

For finite \(X\), remove the last listed premise first and continue
backwards, conditionalizing both question answer sets at every step. By the
one-premise biconditional and induction, the final two sets have the stated
common nested prefix, with an empty basis. Reverse the steps for the converse.
The construction changes both questions; using only a transformed target
would not have transferred the source-answer witnesses just proved.
\end{proof}
\end{annexitem}

\begin{annexitem}{annex:p18.ei.conjunct-syntax-boundary}{thm:p18-ei-conjunct-syntax-boundary}{An explicit language witnessing the syntax limitation}
\begin{proof}
Take the ordinary, uniquely parsed Boolean formula language, without identifying
truth-equivalent formulas as expressions. Let the question category consist
only of \(Q_0\) and questions \(Q_C\) for each declarative \(C\), with
\(dQ_0=\{p,q\}\) and \(dQ_C=\{C\land p,C\land q\}\).
Use Boolean valuations for admissible partitions, so the conjunction truth
clause holds. All assigned answer sets have two syntactically distinct formulas.

The base question is not added-conjunct: its answers are atomic, not
conjunctions with a common prefix. Every other question is added-conjunct
from that base. Thus the restricted CQ1-conjunction condition is met: every
declarative conjunct can be added to the only non-added-conjunct question,
yielding the assigned \(Q_C\). For inverse CQ2-conjunction, the atomic base
answer set cannot have a conjunction prefix. Any prefixed answer set in the
language has, by unique parsing, exactly its displayed prefix \(C\) and
the residual set \(\{p,q\}\). The residual question \(Q_0\) exists, so
inverse closure holds as well.

The proposed nested answers have right-hand immediate subformulas
\(b\land p,b\land q\), not \(p,q\). They are neither the base answer
set nor any \(\{C\land p,C\land q\}\), even after exchanging order.
Hence their question is absent. A bundled question with prefix \(a\land b\)
does exist, but its answers are \((a\land b)\land p\) and
\((a\land b)\land q\), different expressions. Their Boolean equivalence
to the nested expressions does not alter the stipulated question category.
\end{proof}
\end{annexitem}

\begin{annexitem}{annex:p18.ei.conjunct-reduction}{thm:p18-ei-conjunct-reduction}{Conjunctive reduction without an unstated syntax closure}
\begin{proof}
Induction on the finite prefix length and the conjunction truth clause shows
\(G_X(A)\in T_P\) iff \(X\subseteq T_P\) and \(A\in T_P\).
For any \(Y\ne\varnothing\), it follows that
\[
 X,D\mathrel{\|=}_L Y
 \quad\Longleftrightarrow\quad
 \{G_X(D)\}\mathrel{\|=}_L G_X(Y).
\]
Indeed a partition satisfying the premise on the right satisfies exactly
\(X,D\), and an image answer is true there exactly when its unprefixed
answer is true. This proves both directions, not only preservation.

Apply the equivalence with \(D\) each original source answer and \(Y=dR\)
for the first clause. For the second clause apply it with \(D\) a target
answer and \(Y\) its original witness. Injectivity transfers nonemptiness
and properness. Conversely take the inverse image of any transformed proper
witness and apply the reverse equivalence. This proves the result for the
explicitly assumed nested questions.

For the bundled variant, the conjunction truth clause gives
\(C_X\in T_P\) iff \(X\subseteq T_P\). The same proof with the map
\(A\mapsto C_X\land A\) applies. A single addition to each original
non-added-conjunct question is within the source's restricted closure
hypothesis. Iterating that restricted hypothesis on already prefixed questions
is not used. Truth equivalence of the two prefix forms is proved by their
truth conditions; it is not substituted for existence of either question.
\end{proof}
\end{annexitem}

\begin{annexitem}{annex:p18.ei.composition}{thm:p18-ei-composition}{The sound composition routes and their exact witness condition}
\begin{proof}
Let a partition satisfy \(X\cup Z\) and an original answer to \(Q\).
The first leg gives a true answer to \(R\), and the second gives a true
answer to \(S\). Thus the composed first clause holds.

For a direct answer \(C\) to \(S\), regularity of the second leg supplies
\(B\in dR\) with \(Z,C\models_L B\). The first leg supplies a nonempty
proper \(Y_B\subsetneq dQ\) entailed by \(X,B\). At every partition
satisfying \(X,Z,C\), \(B\) is true and then some answer of \(Y_B\) is
true. This is the composed second clause. If the first leg is also regular,
replace \(Y_B\) by its entailed singleton, proving regular composition.

For the more general explicitly supplied witness families, a partition
satisfying \(X,Z,C\) has a true \(B\in V_C\), then a true member of
\(Y_B\). Thus it has a true member of the union. The union is nonempty
because \(V_C\) and all its indexed witnesses are nonempty; it is proper by
the stated additional condition. The same first-clause proof and this union
give implication. No assumption that arbitrary unions of proper subsets
remain proper occurs, nor is an unprovided choice function for an infinite
family silently imported: the witness families are specified input data.
\end{proof}
\end{annexitem}

\begin{annexitem}{annex:p18.ei.nontransitivity}{thm:p18-ei-nontransitivity}{A satisfiable-answer countermodel to ordinary transitivity}
\begin{proof}
Each answer \(p,q,r\) entails some answer of \(R\). The witnesses for
\(p\lor q,p\lor r,q\lor r\) are respectively
\(\{p,q\},\{p,r\},\{q,r\}\), each nonempty and proper. This proves
\(\operatorname{Im}(Q,X,R)\).

Every answer to \(R\) entails \(p\lor q\lor r\), a target answer to
\(S\). For this answer the set \(\{p\lor q,p\lor r\}\) is a nonempty
proper witness in \(dR\), since its disjunction is equivalent to
\(p\lor q\lor r\). For the target answer \(s\), the basis \(s\to p\)
gives \(p\), hence the singleton witness \(\{p\lor q\}\). This proves
\(\operatorname{Im}(R,X,S)\).

To refute \(\operatorname{Im}(Q,X,S)\), consider its answer
\(p\lor q\lor r\). For any proposed proper subset of \(\{p,q,r\}\),
choose an excluded variable to be the only true one, and set \(s=0\).
The basis and target answer are true but every member of the proposed
witness is false. Thus no allowed witness exists. The first relation's
witnesses for the two members of the second relation's witness have union
\(\{p,q,r\}\), exhibiting precisely the lost properness.
All original and target formulas are individually satisfiable. In particular
\(s=1,p=1\) is compatible with the basis, so its answer is not a vacuous
contradiction introduced to manufacture the failure.
\end{proof}
\end{annexitem}

\begin{annexitem}{annex:p18.ei.native-patterns}{thm:p18-ei-native-patterns}{Both native pattern directions and their regularity limits}
\begin{proof}
For the refinement direction, \(p\) supplies itself and \(q\lor r\)
supplies at least one of \(q,r\), giving the first clause. Each target
answer entails a single original answer: \(p\) entails \(p\), while
\(q\) and \(r\) each entail \(q\lor r\). The singleton witnesses
are proper in the two-element original answer set, so implication is regular.

In the coarsening direction, each of \(p,q,r\) entails one coarse answer.
The witnesses for \(p\) and \(q\lor r\) are respectively \(\{p\}\)
and \(\{q,r\}\), nonempty proper subsets of the three-element set.
This proves ordinary implication. To refute regularity for \(q\lor r\),
the assignment \((p,q,r)=(0,1,0)\) refutes entailment of \(p\) and
\(r\), while \((0,0,1)\) refutes entailment of \(q\). Hence this
answer entails no individual source answer, although it entails their stated
two-answer subset.

For auxiliary equivalence, at any partition satisfying \(X,A_i\), the
Boolean biconditional truth clause gives \(B_i\), hence coverage of the
target question. Conversely \(X,B_i\) gives \(A_i\), a singleton proper
witness. This proves regular implication. Swapping the two lists proves
the reverse regular implication by the identical two truth-clause steps.
All lists are well formed by hypothesis, so no duplicates are silently
removed. The two false-basis valuations in the reader make respectively
only \(p\) and only \(q\) true; they falsify the relevant biconditional
and yield the stated opposite soundness mismatches.
\end{proof}
\end{annexitem}

\begin{annexitem}{annex:p18.ei.inquiry}{thm:p18-ei-inquiry}{Every branch of the source's travel pattern}
\begin{proof}
If a partition satisfies \(X\) and any of \(p,l,m\), their disjunction is
true and the first auxiliary premise gives \(a\lor e\). If it satisfies
\(X,a\), the second premise gives an answer in \(\{p,l\}\); if it
satisfies \(X,e\), the third gives the answer \(m\). These are exactly
the coverage and two proper witnesses for \(\operatorname{Im}(Q,X,R)\).
Accepting a true evening answer therefore entails \(m\).

After accepting \(a\), \(Z\) entails \(p\lor l\). The assignments
\((p,l,m,a,e)=(1,0,0,1,0)\) and \((0,1,0,1,0)\) satisfy \(Z\) and
respectively falsify \(l\) and \(p\). Thus the branch question is evoked.
Every Boolean partition supplies an answer to the polar question
\(?\{p,\neg p\}\). Its positive answer entails \(p\), while its negative
answer and \(Z\) entail \(l\) by the truth clause for disjunction. Hence
the branch regularly implies this polar question. The profile theorem also
shows it is evoked before either response is received. On the true response
the new basis entails one direct answer to the original question, so this
branch actually resolves a requested answer.

Omit the second auxiliary premise. The valuation
\((p,l,m,a,e)=(0,0,0,1,0)\) satisfies the remaining basis and the target
answer \(a\), but no original answer. It refutes every possible reduction
witness for \(a\). Thus the weakened inquiry lacks implication, although
an unrelated check that \(a\lor e\) can hold would miss this failure.
Finally, the full premises allow \(p=m=a=1,e=l=0\); the valid morning
witness does not imply \(\neg m\). A unique destination would need a
separate exclusivity premise.
\end{proof}
\end{annexitem}

\begin{annexitem}{annex:p18.ei.wh-cases}{thm:p18-ei-wh-cases}{Wh-answer mapping and concrete domain countercases}
\begin{proof}
The instance-substitution fact used here follows directly by structural
induction on the first-order formula. For an atomic predicate or equality,
replacing a free occurrence of \(x\) by \(c\) has exactly the effect of
assigning \(x\) the interpretation of \(c\); there are no function-symbol
term cases. Boolean constructors preserve this identity by their truth
clauses. Under a quantifier binding \(x\), substitution does not act on
those bound occurrences, and both evaluations assign the quantified domain
individual to \(x\). Under a quantifier binding another variable, evaluate
each domain value and apply the induction hypothesis to the subformula.
Universal and existential truth therefore preserve the identity as well.
Constants introduce no variable capture. Consequently equal denotations of
two constants give the same truth value for their instances of \(Ax\).

In any nonempty-domain model satisfying \(\forall x(Ax\leftrightarrow Bx)\),
the closed instances \(A(x/c)\), \(B(x/c)\) have the same truth value for
each constant \(c\). A true source instance therefore gives its target
instance; a true target instance gives the corresponding original direct
answer, a singleton proper witness. This proves regular implication of
existential wh-questions without asserting that either has a true instance.

For a complete-list answer using constants \(c_1,\ldots,c_k\), the same
universal equivalence transfers every positive instance. It also transfers
the exhaustion sentence: every individual satisfying \(B\) satisfies \(A\)
and therefore equals one of the listed denotations, and conversely. Hence
the entire target list formula is equivalent on the basis to the source list
formula with the same constants. This proves coverage and regular singleton
reduction in both directions. Distinct constant symbols do not need distinct
values for this argument.

For the finite restriction basis, let any \(A(x/c)\) be true. Its denotation
satisfies \(A\), so the universal restriction makes that denotation equal
to a denotation of some \(c_i\). Interpretation of equality and substitution
then give \(A(x/c_i)\). Thus the finite target set is covered. Conversely
each target instance is itself a direct answer of the original wh-question,
giving regularity. These are precisely the required two clauses.

For the unnamed-witness failure take domain \(\{0,1\}\), interpret every
constant as 0, and give predicate \(A\) extension \(\{1\}\). Then
\(\exists xAx\) is true but every named instance is false: the existential
wh-question is unsound. This model uses a nonempty domain and satisfies
ordinary equality. Naming all domain members would exclude this particular
countercase, but is an additional restriction.

For the empty-extension failure take a one-element named domain and empty
extension. No positive instance is true, and every nonempty list answer has
a false positive conjunct. No ``there are none'' sentence is among the source
\(U\)-answers. For the infinite-list failure take the natural numbers as
domain, provide constants naming all its members, and let \(A\) hold of
every member. Any finite proposed list misses another member, falsifying its
exhaustion clause. Thus complete-list soundness fails even with full naming.
These are failures of question soundness, not refutations of the earlier
regular mappings: if the paired predicates have the same extension, both
questions fail soundness together.
\end{proof}
\end{annexitem}

\begin{annexitem}{annex:p18.ei.certificate}{thm:p18-ei-certificate}{Finite semantic certificates without a calculus substitution}
\begin{proof}
Every displayed input formula has a finite parse tree and only finitely many
variables occur in the finite inputs. Recursive Boolean evaluation terminates
on each assignment. Each assignment extends to a full-language valuation by
setting all remaining variables false; every full valuation restricts to an
assignment with the same input-formula values. Thus the two assignment tests
are exactly the two semantic implication clauses for the supplied witnesses.
This proves soundness of every passing certificate.

Conversely, if implication holds, each target answer has a nonempty proper
witness subset of the finite source-answer set by definition. For each of
the finitely many target answers select one such subset; the clauses just
proved make its assignment tests pass. If there are \(m\) source answers,
there are \(2^m-2\) candidate subsets and finitely many tuples of them for
the finitely many target answers. Enumerating those tuples and the finite
assignments therefore terminates and finds a passing one exactly when
implication holds. For regular implication restrict candidates to singletons.
The argument concerns semantic validity in this finite Boolean instance;
it neither provides a general first-order decision procedure nor imports
question-premise provability from PMCE.
\end{proof}
\end{annexitem}
